\subsection{The Error Localization Theorem}\label{the-error-localization-theorem}

\paragraph{Scope of this section.}
This section studies \emph{maintenance/localization} complexity, measured by locations to inspect ($E(\cdot)$), not per-instance identification complexity ($W$) from Sections~\ref{sec:framework} and~\ref{sec:witness}. The metrics are related but distinct: $W$ quantifies online class-identification effort, while $E$ quantifies where constraint logic is distributed in implementations.

\begin{definition}[Error location count]
Let $E(\mathcal{O})$ be the number of locations that must be inspected to find all potential violations of a constraint under observation family $\mathcal{O}$.
\end{definition}

\begin{theorem}[Nominal-tag localization]\label{thm:nominal-localization}
$E(\text{nominal-tag}) = O(1)$.
\end{theorem}
\leanmeta{ \LH{ACS7}, \LH{DOM1}, \LH{FORC2}}

\begin{proof}
Under nominal-tag observation, the constraint ``$v$ must be of class $A$'' is satisfied iff $\tau(v) \in \text{subtypes}(A)$. This is determined at a single location: the definition of $\tau(v)$'s class. One location.
\end{proof}

\begin{theorem}[Declared-attribute localization]\label{thm:declared-localization}
$E(\text{attribute-only, declared}) = O(k)$ where $k$ = number of entity classes.
\end{theorem}
\leanmeta{ \LH{ACS3}}

\begin{proof}
With declared attribute sets, the constraint ``$v$ must satisfy attribute $I$'' requires verifying that each class satisfies all attributes in $I$. For $k$ classes, $O(k)$ locations.
\end{proof}

\begin{theorem}[Attribute-only localization]\label{thm:attribute-localization}
$E(\text{attribute-only}) = \Omega(m)$ where $m$ = number of query sites.
\end{theorem}
\leanmeta{ \LH{ACS3}, \LH{ACS9}, \LH{COH1}}

\begin{proof}
Under attribute-only observation, each query site independently checks ``does $v$ have attribute $a$?'' with no centralized declaration. For $m$ query sites, each must be inspected. Lower bound is $\Omega(m)$.
\end{proof}

\begin{corollary}[Strict dominance]\label{cor:strict-dominance}
Nominal-tag observation strictly dominates attribute-only: $E(\text{nominal-tag}) = O(1) < \Omega(m) = E(\text{attribute-only})$ for all $m > 1$.
\end{corollary}
\leanmeta{ \LH{ACS4}, \LH{DOM1}, \LH{FORC2}}

\subsection{The Information Scattering Theorem}\label{the-information-scattering-theorem}

\begin{definition}[Constraint encoding locations]
Let $I(\mathcal{O}, c)$ be the set of locations where constraint $c$ is encoded under observation family $\mathcal{O}$.
\end{definition}

\begin{theorem}[Attribute-only scattering]\label{thm:attribute-scattering}
For attribute-only observation, $|I(\text{attribute-only}, c)| = O(m)$ where $m$ = query sites using constraint $c$.
\end{theorem}
\leanmeta{ \LH{ACS3}}

\begin{proof}
Each attribute query independently encodes the constraint. No shared reference exists. Constraint encodings scale with query sites.
\end{proof}

\begin{theorem}[Nominal-tag centralization]\label{thm:nominal-centralization}
For nominal-tag observation, $|I(\text{nominal-tag}, c)| = O(1)$.
\end{theorem}
\leanmeta{ \LH{ACS6}, \LH{UNIQ1}, \LH{DOM1}}

\begin{proof}
The constraint ``must be of class $A$'' is encoded once in the definition of $A$. All tag checks reference this single definition.
\end{proof}

\begin{corollary}[Maintenance entropy]\label{cor:maintenance-entropy}
Attribute-only observation maximizes maintenance entropy; nominal-tag observation minimizes it.
\end{corollary}
\leanmeta{ \LH{ACS2}, \LH{COH1}, \LH{FORC2}}
